\label{Chapter Introduction}
% \begin{multicols}{2}
\section{Introduction}
    The \textbf{Virtual Youtuber} (VTuber), which began in Japan and is now quite common in the entertainment category of the Youtube platform, is an emerging digital trend as a result of the rapid advancement of technology over the past several years.
    VTubers use a virtual avatar produced with computer graphics and real-time motion capture software to record their movements in order to control their models and transmit them to their audience instead of streaming with their actual face.
    The majority of VTubers are English- and Japanese-speaking YouTubers, with the most well-known ones being the idols in the group Hololive under Cover Corporation, who have millions of subscribers.
    The name Hololive was first applied to Cover's December 2017-released 3D stream distribution software, and then to its first-generation female VTuber agency. 
    
    As a response, we thought it would be fascinating to develop a face detection project that can distinguish more than 55 different Hololive talents.
    Therefore, fans can explore and play about with this model.
    Furthermore, it makes it easy for curious newcomers to locate and identify the idols they are interested in.
    
    The \textbf{YOLOv5} model, one of the most well-known object identification algorithms because of its speed and accuracy, was chosen for this project.
    After this introduction, we'll go into more depth about the training process, the outcomes, and eventually some discussion.